\setchapter{3}{Ordered Data and Count Data Models}
\ChapterHeading

\section{Ordered data}
In some problems, the variate of interest assumes more than two discrete outcomes, but these are inherently \Emph{orderded}.
\Info{Ordered data is a generalization of the binary choice model but with an order.} 
Zavoina and McElvey (1975) modeled ordered data using the following latent variable framework,
\begin{equation}
    Y^*_i = X'_i\beta_0+u_i, \quad Y_i = \begin{cases}
        0 & Y_i^*\leq \mu_0\\
        1 & \mu_0<Y_i^*\leq\mu_1\\
        2 & \mu_1<Y_i^*\leq\mu_2\\
        \vdots & \vdots\\
        J-1 & \mu_{J-2}<Y_i^*\leq\mu_{J-1}\\
        J & \mu_{J-1}<Y_i^*
    \end{cases}
\end{equation}
where the threshold parameters are such that $0=\mu_0<\mu_1<\cdots<\mu_{J-1}$ and $Y_i^*$ is the latent variable.\\

If \Emph{the distribution of $u_i$ is specified}, the \Emph{unknwon parameters} $\beta$ and $\mu_2,\ldots,\mu_{J-1}$ can be estimated by \Emph{ML}.
Assume that $u_i$ has dsitribution function $F(\cdot)$ and is independent of $X_i$, notice that,
\begin{gather*}
    P_0(X_i,\beta_0) = P(Y_i=0|X_i) = F(-X_i'\beta_0)\\
    P_1(X_i,\beta_0) = P(Y_i=1|X_i) = F(\mu_1-X_i'\beta_0) - F(-X_i'\beta_0)\\
    \vdots\\
    P_j(X_i,\beta_0) = P(Y_i=j|X_i) = F(\mu_j-X_i'\beta_0) - F(\mu_{j-1}-X_i'\beta_0)\\
    \vdots\\
    P_J(X_i,\beta_0) = P(Y_i=J|X_i) = 1 - F(\mu_{J-1}-X_i'\beta_0)
\end{gather*}
Therefore the log-likelihood function is,
\begin{equation*}
    \log L(\theta) = \sum_{i=1}^{n}\sum_{i=1}^{J}1(Y_i=j)\log(P_j(X_i,\beta))
\end{equation*}
As in all discrete choice models, \Emph{the variance of $u_i$ is not identified}.
And the \KeyIdea{ordered-probit and ordered-logit} are the most used special cases of the model.\\

Interpreting coefficients requires some care and is messy. But, the OP and OL models allow us to obtain the \Emph{sign of the partial effects} on $P(Y>j|X_i)$.
For a continuous variable $x_h$ for the OP model,
\begin{equation*}
    \frac{\partial P(Y_i>j|X_i)}{\partial x_h} = \beta_h\Phi(\mu_j - X_i'\beta)
\end{equation*}
If \Emph{$\beta_h>0$, an increase in $x_h$ increases the probability that $Y_i$ is greater than any value j}. Of course, we can interpret the sign of the parameters in the latent variable model.
\Info{It is actually easier to interpret the $\beta$'s of the latent equation to avoid the complications of finding the marginal effect.}\\

A closely related model can be used for \KeyIdea{grouped data}. In this case, the threshold parameters are the limits of the intervals.
The main difference is that for $J>0$, \Emph{the variance of $u_i$ is identified} because the thresholds give information on the scale of $u_i$.
\Info{Only need to estimate the $\beta$'s becasue the $\mu$'s are known.}

\section{Count data models}
In many applications, the variate of interest is the count of the number of occurences of some event in a given period of time.
These data have some very specific characteristics: discreteness, non-negative, and many zeros and a long right-hand tail.

In those conditions, the \Emph{standard linear models are not appropriate} because:
\begin{itemize}
    \item The conditional expectation is necessarily non-negative. \Info{The mean is always positive.}
    \item The data is intrinsically heteroskedastic.
    \item Do not allow the computation of the probability of events of interest.
\end{itemize}

\subsection{The Poisson Regression Model}
The \Term{Poisson regression model} is defined by,
\begin{equation}
    P(Y_i = j|X_i) = \frac{\exp\{-\lambda(X_i.\beta_0)\}\lambda(X_i,\beta_0)^j}{j!},\quad j=0,1,2,\ldots
\end{equation}
with
\begin{gather*}
    E[Y_i|X_i] = \Var(Y_i|X_i) = \lambda(X_i,\beta_0)\quad \textcolor{red}{\text{equidispersion assumption}}\\
    \Var(Y_i) = E_x[\lambda(X_i,\beta_0)] + \Var[\lambda(X_i,\beta_0)]\geq E_x[\lambda(X_i,\beta_0)] = E[Y_i]
\end{gather*}
where in general the specification of $\lambda(X_i,\beta_0) = \exp(X'_i\beta_0)$ is adopted. 
\Info{Because $\exp(X'_i\beta_0)>0$.}
Therefore, 
\begin{equation*}
    \frac{\partial E[Y_i|X_i]}{\partial X_i} = \exp(X'_i\beta_0)\beta_0
\end{equation*}
and \Emph{ML} estimation of $\beta_0$ is straightforward.
\begin{proof}
    Proof of $\Var(Y_i)$
    \begin{align*}
        \Var(Y_i) & \quad\text{using law of total variance}\\
        & = \overbrace{\Var_{X_i}(\underbrace{E[Y_i|X_i]}_{\lambda(\cdot)})}^{\geq 0}
        + E_{X_i}[\underbrace{E[Y_i|X_i]}_{\lambda(\cdot)}]\Rightarrow\\
        \Var(Y_i) & \geq E_{X_i}[\Var(Y_i|X_i)] = E_{X_i}[\lambda(X_i,\beta_0)]\\
        & = E_{X_i}[E[Y_i|X_i]]\quad\text{using the law of interated expectations}\\
        & = E[Y_i]
    \end{align*}
\end{proof}
The log-likelihood function is given by,
\begin{equation*}
    \log L(\beta) = \sum_{i=1}^{n}[-\exp(X'_i\beta)+(X'_i\beta)Y_i-\log(Y_i!)]
\end{equation*}
The \Emph{Hessian is negative definite} for all $X$ and $\beta$ which facilitates the estimation and ensures the unqiueness of the maximum if it exists.
The \Emph{MLE} has the usual properties,
\begin{equation*}
    \sqrt{n}(\hat{\beta}_{ML}-\beta_0)\xrightarrow{d}\to\mcN(0, E[\exp(X'_i\beta_0)X_i X'_i]^{-1})
\end{equation*}
And as usual, \Emph{inference} can be performed using the \KeyIdea{LR, W, and LM tests}.

\subsection{Overdispersion}
The Poisson model imposes equidispersion which is very restrictive.
\begin{definition}
    \textbf{Equidispersion} means the conditional expectation is equal to the variance.\\
    $E[Y_i|X_i] = \Var(Y_i|X_i)$\\
    While \textbf{overdispersion} means that the variance is bigger than the conditional mean.\\
    $\Var(Y_i|X_i)>E[Y_i|X_i]$
\end{definition}
There are many causes for overdispersion: measurement error, misspecification of the condition mean, neglected heterogeneity (random parameter variation).
Focusing on the neglected heterogeneity issue, assuming
\begin{align*}
    E[Y_i|X_i,\epsilon_i] = \exp(X'_i\beta_0+\epsilon_i)&\quad \textcolor{red}{\epsilon_i\text{ can be any possibilities of error}}\\
    E[\exp(\epsilon_i)|X_i] = 1\\
    \Var(\exp(\epsilon_i)|X_i) = \sigma^2
\end{align*}
In this case,
\begin{equation*}
    E[Y_i|X_i] = E[\lambda(X_i,\beta_0)|X_i] = E_{\epsilon}[\exp(X'_i\beta_0+\epsilon_i)|X_i] = \exp(X'_i\beta_0)
\end{equation*}
Therefore, this sort of neglected heterogeneity \Emph{does not change the form of the conditional expectation of $Y_i$}.
\begin{theorem}
    \textbf{Gourieroux, Monfort and Trognon (1984)}\\
    If \Emph{$E[Y_i|X_i] = \lambda(X'_i,\beta_0)$ is correctly specified} and the likelihood function is constructed
    using a probability distribution which does not necessarily correspond to the true distribution of the data,
    but \Emph{belongs to the family of linear exponential distribution}, 
    then the \KeyIdea{Quasi-Maximum Likelihood} estimator is \Emph{consistent} for $\beta_0$.
\end{theorem}
\Info{Wrong estimator $\to$ same exponential family $\to$ still consistent.}
The family of linear exponential distribution includes
\begin{itemize}
    \item Poisson distribution
    \item Normal distribution with fixed variance \Info{ $\to$ using contiuous distribution to describe a discrete distribution.}
    \item Binomial distribution with fixed number of trials
    \item Gamma distribution with fixed shape parameter
\end{itemize}

\Emph{Inference} is done using the results presented for the Quasi-Maximum Likelihood estimator.
In particular, since the Poisson pseudo-MLE is consistent in presence of this sort of misspecification,
valid inference can be based on
\begin{equation*}
    \sqrt{n}(\hat{\beta}-\beta_0)\xrightarrow{d}\mcN(0, A^{-1}BA^{-1})
\end{equation*}
where
\begin{equation*}
    A = E[\exp(X'_i\beta_0X_i X_i')]\quad B = E[(Y_i-\exp(X'_i\beta_0))^2X_i X'_i]
\end{equation*}

Notice that
\begin{align*}
    Var(Y_i|X_i) & = E_{\epsilon_i}[\Var(Y_i|X_i,\epsilon_i)] + \Var(E_{\epsilon_i}[Y_i|X_i,\epsilon_i])\quad \text{law of total variance}\\
    & = \exp(X'_i\beta_0) + \sigma^2\exp(2X'_i\beta_0)\\
    & \geq E[Y_i|X_i] = \exp(X'_i\beta_0)
\end{align*}
This means that the \Emph{presence of overdispersion or underdispersion} can be tested with \KeyIdea{LM test statistic} 
\begin{flalign*}
    & H_0: \sigma^2 = 0 \to \text{equidispersion}\\
    & T = \sum_{i=1}^{n}\frac{(Y_i-\exp(X'_i\hat{\beta}))^2-Y_i}{\sqrt{2\sum_{i=1}^{n}\exp(2X_i'\hat{\beta})}}\xrightarrow{d}\mcN(0,1)
\end{flalign*}
Or a \KeyIdea{RESET test} to test if the specificiation $E[Y_i|X_i] = \exp(X'_i\beta_0)$ is correct.
\begin{flalign*}
    & H_o : \gamma_0 = 0 \to \text{correctly specified thus equidispersion}\\
    & E[Y_i|X_i] = \exp(X'_i\beta_0 + \gamma_0(X'_i\beta_0)^2)
\end{flalign*}

\subsection{Heterogeneity and the Negative Binomial Regression Model}
The assumption that $Y_i$ has a Poisson distribution coniditonal of $X_i$ and $\epsilon_i$ with mean $\lambda_i = \exp(X'_i\beta_0+\epsilon_i)$ leads to the \Term{compound Poisson regression model}.
\begin{flalign}
    & P(Y_i=j|X_i,\epsilon_i) = \frac{\exp[-\exp(X'_i\beta_0+\epsilon_i)]\exp(X'_i\beta_0+\epsilon_i)^j}{j!}\\
    & P(Y_i=j|X_i) = \int_{-\infty}^{+\infty}\frac{\exp[-\exp(X'_i\beta_0+\epsilon_i)]\exp(X'_i\beta_0+\epsilon_i)^j}{j!}g(\epsilon_i)\, d\epsilon_i
\end{flalign}
where $g(\epsilon_i)$ is the density function of $\epsilon_i$ and we assume that $X_i$ and $\epsilon_i$ are independent.
This model can be made operational in different ways:
\begin{itemize}
    \item Pseudo maximum likelihood estimation \Info{$\to$ Refer to previous section.}
    \item Parametric estimation for specified $g(\epsilon_i)$
    \item Semiparametric estimation of $\beta_0$ and $g(\epsilon_i)$
\end{itemize}

Now focusing on the second method when \KeyIdea{$g(\epsilon_i)$ is specified}. \Emph{MLE} can be obtained, btu the estimator \Emph{may not be robust} to departures form the additional distributional assumptions.
Assume that $\exp(\epsilon_i)\sim\Gamma(\sigma^{-2},\sigma^2)$, then $P(Y_i=j|X_i)$ is given by the \Term{Negative-Binomial model}, Cameron and Trivedi (1986).
The Poisson model is obtained as a limiting case when $\sigma^2\to 0$, but $H_0:\sigma^2 = 0$ cannot be tested with a standard LR or Wald test.\\

If the model is \Emph{misspecified} but $E[Y_i|X_i] = \exp(X'_i\beta_0)$ is correct and $\sigma^{-2}$ is fixed, the \Emph{negative-binomial Pseudo-MLE} estimator is \Emph{consistent} for $\beta_0$.
This follows the result from \Term{Gourieroux, Monfort and Trognon (1984)} and the fact that the negative-binomial distribution with $\sigma^{-2}$ fixed is a member of the family of linear exponential distributions.
\begin{lemma} Building on top of the original \textbf{Gourieroux, Monfort and Trognon (1984)}\\
    If the model
    \begin{itemize}
        \item is correct $\xrightarrow{MLE}$ consistent
        \item is incorrect but fix $\sigma^{-2}\xrightarrow{MLE} \hat{\beta}$ consistent
        \item is incorrect $\xrightarrow{MLE} \hat{\beta}, \sigma^2$ inconsistent
    \end{itemize}
\end{lemma}
The \KeyIdea{score test} for $H_0:\sigma^2=0$ is the overdispersion test discussed previously. Other parametric alternatives to the Poisson regression are available.\\

Now focusing on the third method with \KeyIdea{semiparametric estimation} which assumes that $\epsilon$ has a discrete dsitribution with $Q$ support points 
$\alpha_1,\ldots,\alpha_Q$ and corresponding possibilities $\pi_1,\ldots,\pi_Q$, leading to 
\begin{equation}
    P(Y_i=j|X_i) = \sum_{q=1}^{Q}\frac{\exp[-\exp(X'_i\beta_0+\alpha_q)]\exp(X'_i\beta_0+\alpha_q)^j}{j!}\pi_q
\end{equation}
For a given $Q$, estimation of $\beta, \alpha_1,\ldots,\alpha_Q,\pi_1,\ldots,\pi_{Q-1}$ can be performed by \Emph{ML}.
This leads to a \Emph{consistent} estimator if $Q$ is allowed to increase at an appropriate rate.
In practice, the value of $Q$ has to be chosen, for example using information criterion.
But \Emph{inferece is complciated} by the fact that the number of parameters is not fixed.
\Info{This method is nice because the big integral in (35) now becomes $\sum_{q=1}^{Q}\pi_q = 1$ since $\epsilon_i$ is discrete.}

\subsection{Hurdle and Zero-Inflated Poisson Model}
In some cases, the population may be contaminated by individuals for which $Y_i\equiv 0$. There are two ways to model this type of data 
\begin{enumerate}
    \item Zero-Inflated Poisson Model
    \item Hurdle Model
\end{enumerate}

In the \Term{Zero-Inflated Poisson Model}, the zero outcomec can arise fomr one of two regimes. 
In one regime, the outcome is always zero. In the other, the ususal Poisson process is at work.
Let $Z_i$ be a Bernoulli random variable such that 
\begin{equation*}
    Z_i = \begin{cases}
        0 & \text{with }P(Z_i=0|X_i) = p_1\\
        1 & \text{with }P(Z_i=1|X_i) = 1-p_i
    \end{cases}
\end{equation*}
Let $P(Y_i=j|X_i,Z_i=1) = \pi_i(j;\beta_0),\, j=0,1,\ldots$ be the Poisson probability function. And let $P(Y_i=0|X_i,Z_i=0) = 1$.
Note that,
\begin{align*}
    P(Y_i=0|X_i) & \quad \text{law of total probability}\\
    & = P(Z_i=0|X_i)P(Y_i=0|X_i,Z_i=0) + P(Z_i=1|X_i)P(Y_i=0|X_i,Z_i=1)\\
    & = P(Z_i=0|X_i) + P(Z_i=1|X_i)P(Y_i=0|X_i,Z_i=1)\\
    & = p_i + (1-p_i)\pi_i(0;\beta_0)
\end{align*}
For $j>0$,
\begin{align*}
    P(Y_i=j|X_i) & \quad \text{law of total probability}\\
    & = \underbrace{P(Z_i=0|X_i)}_{=p_i}\underbrace{P(Y_i = j|X_i,Z_i = 0)}_{=0 \text{ because all the mass are at 0 so cannot be j}}
    + \underbrace{P(Z_i=1|X_i)}_{=(1-p_i)}\underbrace{P(Y_i=j|X_i,Z_i=1)}_{\pi(j;\beta_0)}\\
    & = (1-p_i)\pi(j;\beta_0)
\end{align*}
Also notice that,
\begin{align*}
    E[Y_i|X_i] & = \sum_{j=0}^{\infty}jP(Y_i=j|X_i) = 0\times P(Y_i=0|X_i) + \sum_{j=1}^{\infty}jP(Y_i=j|X_i)\\
    & = \sum_{j=1}^{\infty}j(1-p_i)\pi(j;\beta_0)\\
    & = (1-p_i)\underbrace{\sum_{j=1}^{\infty}j\pi(j;\beta_0)}_{\exp(X'_i\beta_0)\text{ mean of the Poisson random variable}}
\end{align*}
Therefore, the standard \Emph{pseudo maximum likelihood does not hold here if $p_i$ depends on $X_i$}.
The log-likelihood function this model is written as 
\begin{equation}
    \log L(\beta) = \sum_{i=1}^{n}\log\{[p_i+(1-p_i)\pi_i(0;\beta)]^{1(Y_i=0)}\times[(1-p_i)\pi_i(Y_i;\beta)]^{1(Y_i>0)}\}
\end{equation}

The \Term{Hurdle Model} is a different extention of the basic count data model is obtained by letting the zero and positive observations be generated by different mechanisms.
A binary probability model determines whether a zero or a nonzero outcome occurs, then in the latter case, we observe always a positive integer $1,2,\ldots$
Consider a Bernoulli random variable 
\begin{equation*}
    W_i = \begin{cases}
        1 & \text{with }P(W_i=1|X_i) = 1-q_i\\
        0 & \text{with }P(Z_i=0|X_i) = q_i
    \end{cases}
\end{equation*}
where \Emph{$q_i$ may depend on $X_i$}.
\begin{flalign*}
    & P(Y_i=0|X_i, W_i=0) = 1\\
    & P(Y_i=0|X_i, W_i=1) = 0\\
    & P(Y_i=j|X_i, W_i=1) = \pi^*(j;\beta_0)\quad j=1,2,\ldots
\end{flalign*}
\Info{As the model only takes strictly positive values, $\pi^*(j;\beta_0)$ cannot be Poisson.} In this case,
\begin{align*}
    P(Y_i=0|X_i) & = P(W_i=0|X_i)P(Y_i=0|X_i,W_i=0)+P(W_i=1|X_i)P(Y_i=0|X_i,W_i=1)\\
    & = q_i\times 1 + (1-q_i)\times 0 = q_i
\end{align*}
For $j=1,2,\ldots$
\begin{align*}
    P(Y_i=j|X_i) & = P(W_i=0|X_i)P(Y_i=j|X_i,W_i=0)+P(W_i=1|X_i)P(Y_i=j|X_i,W_i=1)\\
    & = P(W_i=1|X_i)P(Y_i=j|X_i,W_i=1)\\
    & = q_i\times 0 + (1-q_i)\pi^*(j;\beta_0)
\end{align*}
In this case, we have 
\begin{align*}
    E[Y_i|X_i] & = \sum_{j=0}^{\infty}jP(Y_i=j|X_i) = 0\times P(Y_i=0|X_i) + \sum_{j=1}^{\infty}jP(Y_i=j|X_i)\\
    & = \sum_{j=1}^{\infty}j(1-p_i)\pi^*(j;\beta_0)\\
    & = (1-p_i)\underbrace{\sum_{j=1}^{\infty}j\pi^*(j;\beta_0)}_{\text{mean of } \pi^*(j;\beta_0)\text{ cannot be }\exp(X'_i\beta_0)}
\end{align*}
Again, the \Emph{standard pseudo maximun likelihood result does not hold here.} Then, the log-likelihood function has the form
\begin{equation}
    \log L(\beta) = \sum_{i=1}^{n}\{1(Y_i=0)(\log q_i)+1(Y_i>0)\log(1-q_i)+1(Y_i>0)\log[\pi^*(Y_i;\beta)]\}
\end{equation}
Notice that this function is separable meaning the \Emph{correlated unobserved heterogeneity can be allowed for} and integrated out numerically.
\Info{Can model $q_i$'s separately, $q_i = \frac{\exp(X_i'\gamma_0)}{1+\exp(X'_i\gamma_0)}$ with this then we can estiamte the $\beta$'s.}\\

Usually, $\pi^*(j;\beta_0)$ is specified as a truncated Poisson of the form
\begin{equation*}
    \pi^*(j;\beta_0) = \frac{\exp(-\lambda_i)\lambda_i^j}{(1-\exp(-\lambda_i))j!},\quad j>0
\end{equation*}
with $\lambda_i = \exp(X'_i\beta_0)$. However, there is no real truncation in this model, so an equally valid specification is 
\begin{equation*}
    \pi^*(j;\beta_0) = \frac{\exp(-\lambda_i)\lambda_i^{j-1}}{(j-1)!},\quad j>0
\end{equation*}
When the \Emph{truncated Poisson specification is used}, we can test the \KeyIdea{null of no hurdle} through
\begin{flalign*}
    & H_0: \beta_0 = \gamma_0\\
    & q_i = \exp(-\exp(X'_i\gamma_0))
\end{flalign*}
In any case, \Emph{consistency depends on the distributional assumptions.}


% \begin{definition}
% Theorem
% \begin{lemma}
% \begin{Example}
% \begin{proof}
%\newcommand{\KeyIdea}  % 关键想法：浅绿底
%\newcommand{\Term}     % 专业名词：品牌深绿字
%\newcommand{\Emph}     % 强调词汇：红
%\newcommand{\Confuse}    % 不懂/存疑处
%\newcommand{\Info}     % 说明或注解